\begin{sol}[Cuenta deslizante sobre un aro que rota]

\subsubsection*{Pregunta 1: Grados de libertad}

La velocidad angular $\Omega$ del aro es impuesta externamente (por un motor) y no obedece una ecuación dinámica propia: no es una coordenada generalizada, sino un parámetro del vínculo (el sistema es \emph{reónomo}, con dependencia explícita del tiempo a través de la orientación del plano del aro). El único grado de libertad dinámico es la posición de la cuenta sobre el aro, descrita por $\phi$.

\subsubsection*{Pregunta 2: Lagrangiano}

Con el eje de rotación vertical ($z$) y el plano del aro girando con ángulo azimutal $\psi = \Omega t$, la posición de la cuenta en coordenadas cartesianas es:
$$x = R\sin\phi\cos(\Omega t), \qquad y = R\sin\phi\sin(\Omega t), \qquad z = -R\cos\phi$$

Derivando y sumando los cuadrados (los términos cruzados se cancelan, como ocurre en coordenadas esféricas con $r=R$ fijo y $\theta \to \phi$), se obtiene la rapidez al cuadrado:
$$v^2 = R^2\dot\phi^2 + R^2\Omega^2\sin^2\phi$$

Energía cinética:
$$T = \frac{1}{2}mR^2\dot\phi^2 + \frac{1}{2}mR^2\Omega^2\sin^2\phi$$

Energía potencial, con referencia en el centro del aro:
$$V = mgz = -mgR\cos\phi$$

$$\boxed{L = \frac{1}{2}mR^2\dot\phi^2 + \frac{1}{2}mR^2\Omega^2\sin^2\phi + mgR\cos\phi}$$

\subsubsection*{Pregunta 3: Ecuación de movimiento}

$$\frac{\partial L}{\partial\dot\phi} = mR^2\dot\phi \quad\Rightarrow\quad \frac{d}{dt}\left(\frac{\partial L}{\partial\dot\phi}\right) = mR^2\ddot\phi$$

$$\frac{\partial L}{\partial\phi} = mR^2\Omega^2\sin\phi\cos\phi - mgR\sin\phi$$

Euler-Lagrange:
$$mR^2\ddot\phi = mR^2\Omega^2\sin\phi\cos\phi - mgR\sin\phi$$

$$\boxed{\ddot\phi - \Omega^2\sin\phi\cos\phi + \frac{g}{R}\sin\phi = 0}$$

Esta ecuación puede reescribirse como $\ddot\phi = \sin\phi\left(\Omega^2\cos\phi - \dfrac{g}{R}\right)$, revelando que el término centrífugo ($\Omega^2\cos\phi$) compite con el término gravitatorio ($g/R$).

\subsubsection*{Pregunta 4: Equilibrios relativos}

Los equilibrios satisfacen $\ddot\phi = 0$ con $\dot\phi = 0$, es decir:
$$\sin\phi\left(\Omega^2\cos\phi - \frac{g}{R}\right) = 0$$

\textbf{Soluciones triviales:} $\phi = 0$ (fondo del aro) y $\phi = \pi$ (parte superior), que existen para cualquier $\Omega$.

\textbf{Solución no trivial:} $\cos\phi_{eq} = \dfrac{g}{R\Omega^2}$, que tiene solución física ($|\cos\phi_{eq}| \leq 1$) solo si:

$$\boxed{\Omega^2 \geq \Omega_c^2 = \frac{g}{R}}$$

\paragraph{Interpretación física} Para $\Omega < \Omega_c$, la única posición de equilibrio estable es el fondo del aro ($\phi=0$), como es intuitivo. Al superar la velocidad angular crítica $\Omega_c = \sqrt{g/R}$, el efecto centrífugo domina cerca del fondo, el equilibrio en $\phi=0$ se vuelve inestable, y aparece un nuevo equilibrio estable en $\phi_{eq} = \arccos(g/(R\Omega^2))$, que se aleja progresivamente del fondo a medida que $\Omega$ aumenta. Este es un ejemplo clásico de bifurcación (transición de estabilidad) en un sistema mecánico simple.

\end{sol}
